% Irodov Part V §5.7, original problems 246--257.
% Source Arihant PDF printed pp.227--228, answers printed pp.337.
% The original printed problem 254 has omega^2, but answer uses omega^3.
% Both literal and physically corrected interpretations are documented.

\thermalopticsproblem{246}{227}
利用维恩的热辐射普遍标度律
$u_\omega(\omega,T)=\omega^3F(\omega/T)$，
证明：(a) 最概然辐射角频率与 $T$ 成正比；
(b) 单位角频率谱能量密度的最大值与 $T^3$ 成正比；
(c) 黑体总辐出度与 $T^4$ 成正比。
\exsolution
令 $x=\omega/T$，
则 $\omega=xT$。
谱能量密度可写成
\[
 u_\omega=T^3x^3F(x).
\]
对固定温度求峰值，
最大值对应的 $x=x_m$ 是函数 $x^3F(x)$
的温度无关极大点，故
\[
 \boxed{\omega_m=x_mT\propto T},\qquad
 \boxed{(u_\omega)_{\max}\propto T^3}.
\]
总能量密度
\[
 u(T)=\int_0^\infty u_\omega\,d\omega
 =T^4\int_0^\infty x^3F(x)\,dx.
\]
对各向同性平衡辐射，
总辐出度 $M_e=cu/4$，因此
\[
 \boxed{M_e\propto T^4}.
\]
\excheck 此处假设积分收敛，
且形状函数 $F$ 本身不显含温度；无需先采用普朗克公式。

\thermalopticsproblem{247}{227}
两个加热黑体中，一个温度为
$T_1=2500\,\mathrm K$。
另一个黑体的辐出度峰值对应波长比第一个大
$\Delta\lambda=0.50\,\mu\mathrm m$。
求第二个黑体的温度。
\exsolution
维恩位移律
$\lambda_m T=b$，取
$b=2.898\times10^{-3}\,\mathrm{m\,K}$。
由
$b/T_2=b/T_1+\Delta\lambda$，
可得
\[
 \boxed{T_2=
 \frac{bT_1}{b+T_1\Delta\lambda}}.
\]
代入
$\Delta\lambda=5.0\times10^{-7}\,\mathrm m$，
\[
 \boxed{T_2\simeq1.75\times10^3\,\mathrm K}.
\]
\excheck 峰值波长变长，温度应降低；
本题 $T_2<T_1$ 满足这一极限检查。

\thermalopticsproblem{248}{227}
某黑体的总辐出度为
$M_e=3.0\,\mathrm{W/cm^2}$。
求该黑体光谱辐出度最大值对应的波长。
\exsolution
首先换算为
$M_e=3.0\times10^4\,\mathrm{W/m^2}$。
斯特藩--玻尔兹曼定律给
$M_e=\sigma T^4$，所以
$T=(M_e/\sigma)^{1/4}\simeq853\,\mathrm K$。
再由维恩位移律
\[
 \boxed{\lambda_m=
 b\left(\frac{\sigma}{M_e}\right)^{1/4}
 \simeq3.40\,\mu\mathrm m}.
\]
\excheck 用的是对\emph{波长谱}的峰值所对应的
位移常数 $b$，不能改用频率谱峰值的不同常数。

\thermalopticsproblem{249}{227}
太阳的辐射谱可近似看作黑体谱，
其波长谱峰值位于
$\lambda_m=0.48\,\mu\mathrm m$。
估算太阳每秒因辐射失去的质量，
以及按此恒定辐射率失去总质量的 $1\%$
所需时间。
\exsolution
太阳表面有效温度由维恩位移律得到
\[
 T\simeq\frac{b}{0.48\times10^{-6}}
 \simeq6.04\times10^3\,\mathrm K.
\]
取太阳半径
$R_\odot\simeq6.96\times10^8\,\mathrm m$，
光度（总辐射功率）
$L=4\pi R_\odot^2\sigma T^4
\simeq4.6\times10^{26}\,\mathrm W$。
由质能等价
$dM/dt=-L/c^2$，得
\[
 \boxed{\left|\frac{dM}{dt}\right|
 \simeq5.1\times10^9\,\mathrm{kg/s}}.
\]
若以
$M_\odot\simeq1.99\times10^{30}\,\mathrm{kg}$
且忽略太阳光度、半径演化，
对应 $1\%$ 质量损失时间为
\[
 \boxed{\Delta t=
 \frac{0.01M_\odot c^2}{L}
 \simeq1.2\times10^{11}\,\mathrm{yr}}.
\]
\excheck 这是\emph{固定光度的形式外推}；
真实太阳演化不允许把该数值当成寿命预测。

\thermalopticsproblem{250}{227}
完全电离的氢等离子体，密度为
$\rho=0.10\,\mathrm{g/cm^3}$。
求热辐射压与等离子体气体动理论压强相等时
的温度 $T$。
取平衡辐射压 $p_{\rm rad}=u/3$，
其中辐射能量密度 $u=4\sigma T^4/c$；
等离子体可视为理想气体。
\exsolution
完全电离氢每个氢核伴随一个自由电子，
忽略电子质量，每单位体积粒子总数
$n_{\rm tot}=2\rho/m_p$
（或者写 $2\rho N_A/M_H$），于是
\[
 p_{\rm gas}=n_{\rm tot}k_BT
 =\frac{2\rho R_g}{M_H}T.
\]
辐射压
$p_{\rm rad}=4\sigma T^4/(3c)$。
令二者相等，约去温度 $T>0$：
\[
 \frac{4\sigma T^3}{3c}
 =\frac{2\rho R_g}{M_H},
 \qquad
 \boxed{T=\left(
 \frac{3c\rho R_g}{2\sigma M_H}
 \right)^{1/3}}.
\]
代
$\rho=100\,\mathrm{kg/m^3}$、
$M_H\simeq1.008\times10^{-3}\,\mathrm{kg/mol}$，
\[
 \boxed{T\simeq1.9\times10^7\,\mathrm K}.
\]
\excheck 氢的\emph{质子和电子各贡献一份气体压强}，
因子2不可漏掉；核反应与电子简并均未计入。

\thermalopticsproblem{251}{227}
一颗直径 $d=1.2\,\mathrm{cm}$ 的铜球放在
内壁绝对零度且抽成真空的容器中。
初温 $T_0=300\,\mathrm K$，视铜球表面为理想黑体，
求降到初温的 $1/\eta$（$\eta=2.0$）所需时间。
\exsolution
设铜的质量密度为 $\rho$、比热容为 $c_p$，
铜球质量
$m=\rho \pi d^3/6$，
辐射面积 $A=\pi d^2$。
忽略物性温度变化，能量平衡
\[
 mc_p\frac{dT}{dt}=-\sigma AT^4.
\]
分离变量从 $T_0$ 积分到
$T_f=T_0/\eta$：
\[
 \int_{T_f}^{T_0}\frac{dT}{T^4}
 =-\frac{\sigma A}{mc_p}\int_{t}^{0}\!dt,
\]
等价地
$\tfrac13(T_f^{-3}-T_0^{-3})
=\sigma At/(mc_p)$。
所以
\[
 \boxed{t=\frac{c_p\rho d}{18\sigma T_0^3}
 (\eta^3-1)}.
\]
取 $\rho\simeq8.96\times10^3\,\mathrm{kg/m^3}$、
$c_p\simeq385\,\mathrm{J/(kg\,K)}$，
\[
 \boxed{t\simeq1.05\times10^4\,\mathrm s
 \simeq2.9\,\mathrm h\approx3\,\mathrm h}.
\]
\excheck 原题理想黑体与真空壁为0K的重要条件
使净辐射功率恰为 $\sigma AT^4$；
若环境温度不为0，应减去环境辐射项。

\thermalopticsproblem{252}{227--228}
两个黑体腔各有一个直径
$d=1.0\,\mathrm{cm}$ 的小圆孔，
两孔彼此同轴、相向放置，相距
$l=10\,\mathrm{cm}$（原图5.39），
腔外表面理想反射。
第一个腔维持
$T_1=1700\,\mathrm K$。
求第二个腔达到稳态时的温度。
考虑黑体孔的余弦发射规律。
\exsolution
黑体孔的总辐出度为 $\sigma T^4$，
近似平行的小圆孔之间的辐射视角因子为
\[
 F_{1\to2}\simeq\frac{A_2}{\pi l^2}
 =\frac{\pi d^2/4}{\pi l^2}
 =\frac{d^2}{4l^2}.
\]
故第二腔每秒通过孔2吸收的来自孔1功率为
$\sigma T_1^4 A_1 F_{1\to2}$。
稳态时孔2向外发出
$\sigma T_2^4 A_2$，
由于两孔面积相同 $A_1=A_2$，所以
\[
 T_2^4=T_1^4\frac{d^2}{4l^2}.
\]
开四次方得
\[
 \boxed{T_2=T_1\sqrt{\frac d{2l}}
 =1700\sqrt{\frac{1.0}{20}}
 \simeq380\,\mathrm K\approx0.4\,\mathrm{kK}}.
\]
\excheck 应开四次方，故温度是
$T_1\sqrt{d/(2l)}$；
若在书末纯文本中见到 $T_1d/(2l)$，
属于根号在文本提取时丢失。

\thermalopticsproblem{253}{228}
一个容积 $V=1.0\,\mathrm L$ 的黑体腔，
充满温度 $T=1000\,\mathrm K$ 的平衡热辐射。
求 (a) 定容热容 $C_V$；
(b) 辐射的熵 $S$。
\exsolution
平衡辐射能量密度
$u=4\sigma T^4/c$，
故
$U=uV=4\sigma VT^4/c$，
定容热容
\[
 \boxed{C_V=\left(\frac{\partial U}{\partial T}\right)_V
 =\frac{16\sigma VT^3}{c}}.
\]
黑体光子气满足
$p=u/3$，并有
$TS=U+pV=4U/3$（化学势为0），所以
\[
 \boxed{S=\frac{16\sigma VT^3}{3c}}.
\]
以 $V=10^{-3}\,\mathrm{m^3}$ 代入，
\[
 \boxed{C_V\simeq3.0\,\mathrm{nJ/K}},\qquad
 \boxed{S\simeq1.0\,\mathrm{nJ/K}}.
\]
\excheck 熵与热容在本问题分别随
$T^3$ 增长，且 $S=C_V/3$。

\thermalopticsproblem{254}{228}
原书给出如下所谓“维恩谱分布”：
\[
 u_\omega(\omega,T)=A\omega^2e^{-a\omega/T},
 \qquad a=7.64\,\mathrm{ps\,K},
\]
求 $T=2000\,\mathrm K$ 时
(a) 频率谱的最概然\emph{角频率}；
(b) 波长谱的最概然波长。
注意本版题面印出的幂为2，但书末答案按幂为3计算。
\exsolution
\textbf{严格按印刷题面 $\omega^2$ 解。}
设 $D=a/T$，频率谱的极大条件为
\[
 \frac{d}{d\omega}
 (\omega^2e^{-D\omega})=0
 \iff\omega(2-D\omega)=0,
\]
其中 $\omega>0$，故
\[
 \boxed{\omega_{\max,\omega}=\frac{2T}{a}
 \simeq5.24\times10^{14}\,\mathrm{s^{-1}}}.
\]
转换波长谱必须保留雅可比因子：
$\omega=2\pi c/\lambda$、
$|d\omega/d\lambda|=2\pi c/\lambda^2$，
故
\[
 u_\lambda\propto\lambda^{-4}
 e^{-2\pi ca/(T\lambda)}.
\]
对 $\lambda$ 求导，得到
\[
 \boxed{\lambda_{\max,\lambda}=\frac{\pi ca}{2T}
 \simeq1.80\,\mu\mathrm m}.
\]
\textbf{若采用物理上常见的维恩谱
$u_\omega=A\omega^3e^{-a\omega/T}$（与原书答案一致），}
则频率谱与波长谱峰值分别为
\[
 \boxed{\omega_{\max,\omega}^{(3)}
 =\frac{3T}{a}
 \simeq7.85\times10^{14}\,\mathrm{s^{-1}}},
\]
\[
 \boxed{\lambda_{\max,\lambda}^{(3)}
 =\frac{2\pi ca}{5T}
 \simeq1.44\,\mu\mathrm m}.
\]
\excheck \emph{不能}既照印刷题面使用 $\omega^2$，
又照书末抄 $\omega^3$ 的数值。
而且在两种模型中，频率谱峰值和波长谱峰值
不能靠简单的 $\lambda=2\pi c/\omega_{\max}$ 相互转化，
因为谱密度变换引入了雅可比因子。

\thermalopticsproblem{255}{228}
从普朗克黑体辐射公式推导角频率谱能量密度
$u_\omega$ 的两个渐近形式：
(a) $\hbar\omega\ll k_BT$（瑞利--金斯极限）；
(b) $\hbar\omega\gg k_BT$（维恩极限）。
\exsolution
普朗克公式
\[
 u_\omega=
 \frac{\hbar\omega^3}{\pi^2c^3}
 \frac{1}{e^{\hbar\omega/(k_BT)}-1}.
\]
记 $x=\hbar\omega/(k_BT)>0$。
(a) $x\ll1$ 时
$e^x-1=x+O(x^2)$，
因而
\[
 \boxed{u_\omega\simeq
 \frac{k_BT\,\omega^2}{\pi^2c^3}}.
\]
(b) $x\gg1$ 时
$e^x-1\simeq e^x$，
故
\[
 \boxed{u_\omega\simeq
 \frac{\hbar\omega^3}{\pi^2c^3}
 e^{-\hbar\omega/(k_BT)}}.
\]
\excheck 两式都具有单位
$\mathrm{J\,m^{-3}\,(rad/s)^{-1}}$，
且后者的指数必须无量纲。

\thermalopticsproblem{256}{228}
将普朗克公式的角频率谱能量密度
$u_\omega$ 改写为普通频率 $\nu$
以及真空波长 $\lambda$ 的谱能量密度
$u_\nu$、$u_\lambda$。
\exsolution
由于谱元的能量必须相等，
\[
 u_\omega\,d\omega=u_\nu\,d\nu
 =-u_\lambda\,d\lambda
\quad\text{（最后一个等式的负号来自变量反向）}.
\]
普通频率与角频率的关系
$\omega=2\pi\nu$、
$d\omega/d\nu=2\pi$。
令 $h=2\pi\hbar$，得到
\[
 \boxed{u_\nu(\nu,T)=
 \frac{8\pi h\nu^3}{c^3}
 \frac{1}{e^{h\nu/(k_BT)}-1}}.
\]
再以 $\nu=c/\lambda$，
取绝对雅可比因子
$|d\nu/d\lambda|=c/\lambda^2$，
所以
\[
 \boxed{u_\lambda(\lambda,T)=
 \frac{8\pi hc}{\lambda^5}
 \frac{1}{e^{hc/(\lambda k_BT)}-1}.
\]
\excheck $u_\nu$ 与 $u_\lambda$ 的数值单位不同；
例如 $u_\lambda$ 是单位\emph{波长}下的能量密度，
不能不带雅可比因子直接用 $\nu=c/\lambda$ 代入。

\thermalopticsproblem{257}{228}
黑体温度为 $T=3000\,\mathrm K$。
求其\emph{单位面积}在接近
波长谱能量密度最大值处的
$\Delta\lambda=1.0\,\mathrm{nm}$ 窄波段内
所辐射的功率。
\exsolution
波长谱辐出度
$M_\lambda=(c/4)u_\lambda$，
由普朗克公式得
\[
 M_\lambda=
 \frac{2\pi hc^2}{\lambda^5}
 \frac{1}{e^{hc/(\lambda k_BT)}-1}.
\]
波长谱峰值位于
$\lambda_m=b/T$，
取 $b=2.89777\times10^{-3}\,\mathrm{m\,K}$，
有
$\lambda_m\simeq9.659\times10^{-7}\,\mathrm m$。
因为 $\Delta\lambda\ll\lambda_m$，
窄波段辐出度近似为
\[
 \Delta M\simeq M_\lambda(\lambda_m)\Delta\lambda.
\]
将 $T=3000\,\mathrm K$、
$\Delta\lambda=10^{-9}\,\mathrm m$
及精确常数 $h,c,k_B$ 代入，
得
\[
 \boxed{\Delta M\simeq
 3.13\times10^3\,\mathrm{W/m^2}
 =0.313\,\mathrm{W/cm^2}}.
\]
\excheck 相比黑体总辐出度
$\sigma T^4\simeq4.59\times10^6\,\mathrm{W/m^2}$，
这个1nm带宽只占很小部分，数量级合理。
